% from ega4-13-fr.tex
\oldpage[IV-3]{193}
  \item On a l'égalité
  \[
    \dim(\mathcal O_{X',x})=\dim(\mathcal O_{Y',y})
    +\dim(\mathcal O_{X',x}\otimes_{\mathcal O_{Y',y}}k(y))
  \]
  (condition toujours impliquée par (iii) lorsque $\mathcal O_{Y,y}$ est un
  anneau universellement caténaire).
\end{enumerate}
De plus, $X'$ est alors une composante irréductible de $X$ et $Y'$ une
composante irréductible de $Y$.
\end{proposition}

Dans l'énoncé, les anneaux locaux $\mathcal O_{X',x}$ et
$\mathcal O_{Y',y}$ sont relatifs aux sous-préschémas fermés réduits de $X$,
$Y$ respectivement ayant $X'$ et $Y'$ pour espaces sous-jacents respectifs.

En outre, cela prouve que

\begin{corollary}[13.2.8]
\label{IV.13.2.8-fr}
Si les deux membres de (13.2.7.1) sont égaux, les parties fermées irréductibles
$X'$ de $X$ contenant $x$ et vérifiant les conditions (i) à (iv) de (13.2.7)
sont exactement celles pour lesquelles on a
$\dim(\mathcal O_{X',x})=\dim(\mathcal O_{X,x})$ (ce qui implique
nécessairement que $X'$ est une composante irréductible de $X$).
\end{corollary}

La proposition (13.2.7) entraîne~:

\begin{corollary}[13.2.9]
\label{IV.13.2.9-fr}
Soient $Y$ un préschéma localement noethérien, $f:X\to Y$ un morphisme de type
fini, $x$ un point de $X$, $f(x)=y$. Supposons que $\mathcal O_{Y,y}$ soit un
anneau universellement caténaire. Alors les conditions suivantes sont
équivalentes~:
\begin{enumerate}[label=\emph{\alph*)}]
  \item Les deux membres de (13.2.7.1) sont égaux.
  \item Il existe une composante irréductible $Z$ de $f^{-1}(y)$ contenant
  $x$, de dimension $\dim_x(f^{-1}(y))$, telle que, pour tout $x'$ dans un
  voisinage de $x$ dans $Z$, on ait
  \[
  \tag{13.2.9.1}
  \label{IV.13.2.9.1-fr}
    \dim(\mathcal O_{X,x'})=\dim(\mathcal O_{Y,y})
    +\dim(\mathcal O_{X,x'}\otimes_{\mathcal O_{Y,y}}k(y)).
  \]
  \item Il existe une composante irréductible $Z$ de $f^{-1}(y)$ contenant
  $x$, de dimension $\dim_x(f^{-1}(y))$, telle que pour le point générique $z$
  de $Z$, on ait
  \[
  \tag{13.2.9.2}
  \label{IV.13.2.9.2-fr}
    \dim(\mathcal O_{X,z})=\dim(\mathcal O_{Y,y}).
  \]
\end{enumerate}
\end{corollary}

Montrons que $a)$ entraîne $b)$. Soit $\dim_x(f^{-1}(y))=n$~; en vertu de
(13.2.7), il existe une composante irréductible $X'$ de $X$ vérifiant les
conditions (i) à (iv) de (13.2.7)~; soit $Z$ une composante irréductible de
dimension $n$ de $f^{-1}(y)$, contenant $x$ et contenue dans $X'$. Comme
$f^{-1}(y)$ est localement noethérien, il existe un voisinage ouvert $U$ de
$x$ dans $Z$ tel que $U$ ne rencontre aucune composante irréductible de
$f^{-1}(y)$ autre que celles qui contiennent $x$, donc (4.1.1.3)
$\dim_{x'}(f^{-1}(y))=n$ pour tout $x'\in U$~; il est clair alors que les
conditions (i) à (iii) de (13.2.7) sont satisfaites quand on y remplace $x$
par un point quelconque $x'\in U$, et il en est de même de la condition (iv)
puisque $\mathcal O_{Y,y}$ est universellement caténaire~; d'où la conclusion
par (13.2.7). La condition $b)$ entraîne trivialement $c)$ en vertu de
(5.1.2). Enfin, si $c)$ est vérifiée et si $X''$ est une composante
irréductible de $X$ contenant $Z$ et telle que les conditions (i) à (iv) de
(13.2.7) soient vérifiées quand on y remplace $X'$ par $X''$ et $x$ par $z$,
il est clair que ces conditions sont aussi vérifiées pour $X''$ et $x$ puisque
$\mathcal O_{Y,y}$ est universellement caténaire, donc $c)$ implique $a)$.

\begin{proposition}[13.2.10]
\label{IV.13.2.10-fr}
Soient $Y$ un préschéma irréductible localement noethérien, $\eta$ son point
générique, $f:X\to Y$ un morphisme de type fini, $y$ un point de $f(X)$.
Soient $Z_i$ les

